\section{HPQ Review}
\label{research:ReviewHPQ}
Prior to sitting my GCSEs I had the opportunity to complete a higher project qualification(HPQ) in which I focused on the impact of AI and how it could impact our future workforces. With my EPQ ideas leaning towards a similar topic i thought it would be useful to review my previous work, both to see how my have changed over time but also to see if there are any useful insights I can take from it. Please find my prior HPQ report below:
\\
\begin{center}
    \textbf{Will Artificial Intelligence Render Our Workforce Obsolete?}\\
    By Henry Field \autocite{fieldWillArtificialIntelligence}
\end{center}

\sloppy
When looking to the future of our workforce, many media outlets and notable figures will point to a new era of machines and how they will revolutionise how we, as humans will manage our planet and live our day-to-day lives. I have chosen to research this subject to try and decipher what may lie ahead and how this future will impact our lives. Based on our current uptake I think that AI will without doubt cause some major change to global industry but that this change is dictated not by the technology, but by how it is embraced by our population. My argument is based on past historical trends as well as my interpretation of both economical and philosophical forecasts. Firstly, I will introduce the concepts and importance of AI and will then move into my interpretation of recent events.

\subsection{Definitions}
As defined by the Organisation for Economic Co-operation and Development (OECD), AI is “a machine-based system that is capable of influencing Its environment by producing an output for a given set of objectives.” (OECD, 2019) In contrast, Britannica states AI as being “the ability of a digital computer or computer-controlled robot to perform tasks commonly associated with intelligent beings.” (Britannica, 2023) Despite this definition varying between organisations and governments these themes of self-produced influence carry throughout. A worker on the other hand is defined by UK and international law as being a person with a “contract or other arrangement to do work or services in exchange for a reward” (UK Government, 2023). These key themes of self-produced influence and reward are my HPQ focus points. Said to be “the most important technological advancement in decades” by Bill Gates (Gates, 2023), AI has recently made large leaps forward toward bridging the ever-narrowing divide between computational logic and human independent thought. Since the dawn of computing, scientists have been studying the concept of a machine capable of personifying rational responses while learning from its own intuition and it can now be argued that we have reached this milestone. These innovations are predicted by many to hugely impact our society for the better however I have chosen to study the more concerning side of how these Living Machines may influence our ability to gain income and possibly negate the need for human involvement in managing our society.

\subsection{History}
A brainchild of Alan Turing, true AI was first proposed in his 1950 paper, “Computing Machinery and Intelligence” which introduced the now ubiquitous Turing Test. This conceptual examination aimed to form a machine so adapt at reproducing human behaviour that it could communicate with people without their realisation of its own mechanical origins. No functional device was ever built but it laid the foundations for a pathway that would slowly lead us to present-day developments. Independently designed and built in 1950, by researcher Claude Shannon was Theseus, a primitive machine learning mouse that could navigate mazes via trial and error while outputting its “memories” to a telephone relay circuit enabling it to “remember” final solutions. Despite not being true AI, it was revolutionary for the time and due to its ability to improve its own solutions, first entertained the thought that computers could learn from past encounters. Theseus met the machine learning definition as given by IBM as, “a branch of artificial intelligence and computer science focused on the use of data and algorithms to imitate the way humans learn”. ML utilises data storage and algorithmic processing to help a machine refine its output when given similar instruction sets and has previously been the limit of where our technology has been headed towards. Computer scientists such as Arnaud Miribel see AI as having four building blocks: quantity of data, Machine learning, Computing power, and Human expertise (Miribel, 2023) – and that once humanity can combine these aspects, we will have a system sufficiently advanced to be defined as true AI - so let’s break these down:

Data is said to be more valuable than money by some which has caused massive organisations (such as Meta, the parent company of Facebook) to release free products that sustain themselves via the collection and sale of this precious resource. Around 120 zettabytes of data are estimated to be transferred yearly and these statistics are only increasing [fig 1]. Google estimates that around 1.2million data points are required for Ai to function and so I believe it can be safely said that we have developed the capacity to form the large data sets required to train AI models.

Machine learning as previously mentioned has continued to grow as models such as AlphaGo by Google have begun dominating gaming circuits. Despite the rules of Go being relatively straightforward, simply adding stones to a 19x19 grid to gain territory, the average game lasts for 150 moves and cumulates to $10^{170}$ board configurations (google ai study)  more possible results than atoms in the universe. This means the traditional algorithms that chess (a relatively simpler game) "AI" have relied on cannot be used as the optimal move processing time would last many years. These massive achievements showcase the level of technology now available to our workforce and highlights AI's ability to quickly grow and improve.

Our capacity to produce computational power has also massively increased, transistor sizes (a common factor measured for the performance of a processor) which started at 5micrometers may have seemed small but are now bloated when compared to the cutting edge 0.34 nanometres chiplets – a 99.9932% reduction in size.
As the final factor, Human expertise is a rather difficult philosophical principle to tackle as we are still not sure what this entails however this has been simplified by scientists to only include the ways we as the population solve challenges and interact with one another. This includes the processing of “natural language”, which entails how humans would normally converse with one another in different situations (such as written language or verbal conversation) rather than the strictly structured languages such as Python or C which computers have been built to understand. Models that have successfully incorporated these techniques have only recently started to appear in the public domain and are currently facing limited success rates as many companies battle to make their implementation the standard.

Despite these emerging technologies being completely unique to anything humanity has created before, we can still investigate the past to forecast what the future may hold. The structure of the workforce has been constantly shifting since the dawn of humanity, from delegated tasks in primal tribes to the modern-day reality complex corporate structures. AI serves to once again disrupt this by slotting machines inside of human ranks – there hasn’t been such a dramatic shift since the Victorian Era’s world of machines.

\subsection{Prior revolutions}
Seen as one of the earliest accounts of a mechanical boom, the 1760s brought about massive change in Britain as the Industrial Revolution shook the growing populace. Manual labour which had employed large swathes of the workforce was now being replaced by steam and ion, the engine (a widely implemented machine) heavily reduced the reliance on horse-drawn carriage and led to trains now carrying 100 passengers whilst requiring only 3 human operators, heavily reducing ongoing costs in to the (passengers : operators) 5  :1 ratio of traditional methods. Many predicted that these changes would soon destroy the workforce, quickly enveloping every trade but these ideologies never came to fruition due to a singular inability of machinery – the ability to independently think. The engine excelled at its singular job of providing drive however required human operators to direct and fuel it while the loom required human ingenuity to set up and operate. These limitations led to entirely new industries of operations and maintenance – massively reducing the job-seeking population and enabling citizens who are challenged by physical disability to actively partake in the workforce. This shifting of common trades from that of labourers to highly skilled operators allowed these new tools to become aids to amplify the skill and ability of a singular human with increased the nation’s manufacturing capacity greatly. Despite this positive precedent of greater equality and opportunity in the workforce, It has been predicted by Teamstage that 1.7 Million manufacturing jobs have been lost to automation since, with 2.25million industrial robots currently being in use worldwide and that by 2030, 20 million manufacturing jobs could be lost (Teamstage, 2023). As companies have scaled, so have their mechanical operations and these machines have slowly consumed many of the previously accessible jobs.

\subsection{Modern advancements}
In the past ten years, many startups have formed which aimed to focus on improving AI technologies. Open AI is one such company, with their primary product, Chat GPT. This is one of the first examples of the technology that has become freely available to the general population and has shown how powerful AI can be. Currently capable of producing written reports, rewriting essays, and forming complex programs based solely on a written prompt. Despite this apparent success these machines are very much in their beta stages and have a habit of being confidently wrong, Google’s product Bard – for example, has notoriously struggled to answer the question of the first photo taken in space correctly (misidentifying the satellite to be the James Webb Space Telescope rather than the correct, European Southern Observatory’s VLT) despite attempts to rectify this by google. The technology is not yet perfect and it not possible to tell how long It will take for these machines to be reliable enough for constant use, this alone has given many researchers pause when giving timeframes around adoption who see this unpredictable behaviour as being a roadblock to the largescale uptake of AI.

Despite this, many companies are faced with fear of missing out and have rushed to begin working on their implementation of AI within their products and workflows. A recent report by the MIT Task Force on the Work of the Future point to these historical trends as a guide for how the next ten years will play out, a slow but sustainable uptake of AI (MIT - Work of the Future, 2020). However, I believe that we – as a population will see much more explosive growth in human replacement and massive changes in the total global workforce. Choosing to see this trend line of AI adoption as being exponential rather than linear, I instead believe that these past principles are not able to be correctly applied to our modern situation and but should still be used to inform policy around easing the pressure on targeted job sectors due to the much broader capabilities in the implementation of recent advancements. This opinion is backed by researchers such as the CEO of Sinovation Venture, Kai-Fu Lee who wrote that 50% of all current jobs will be automated within 15 years (Lee, 2018)and the World Economic Forum who believe that 97 million jobs centred around AI will be created by 2025 (World Economic Forum, 2020). These recent advancements have quickly pushed past the scope of current legislation which caused many (over 1100) notable figures (politicians and scientists alike) to call for emergency measures to slow down these developments to give time for forward-thinking legislation to be written. Notable signers include Elon Musk, Steve Wozniak and Tristan Harris who have called for a worldwide ban on the advancement of AI for 6 months.

A recent report by Adecco Group (the world largest recruitment agency) found that only 4% of workers believe their job will imminently be replaced by artificial intelligence (Adecco Group, 2022). This statistic concerns many as the number of jobs in which AI could replace humans has been predicted to be as high as 85% from some studies and points to a lack of knowledge from the public, AI could become a massive aid to human workers if they understand how to use it as a tool to further upskill themselves. Many are concerned that this necessity for education will widen the pre-existing knowledge gap which can cause mass inequality. To quote Yann LeCun, “our intelligence is what makes us human, and AI is an extension of that quality” (LeCun, 2022). The UK Government’s report on the impact of AI on jobs (UK Government, 2022) stated that those with only GCSE qualifications will see a 7% reduction in employment in the next 5 years and with recent strike action from numerous unions and calls for higher salaries, employers are looking for AI to fill vacant positions and help decrease running costs. 63% of surveyed organisations in 2017 stated that pressure to reduce costs will require the use of AI and a further 84% said that AI would allow them to obtain / sustain a competitive advantage (Sam Ransbotham, n.d.). [fig 2]

\subsection{My observations}
In conclusion, I believe that the future of AI in the workforce is not all doom and gloom, but that Governments need to cooperate to guide legislation in ensuring a safe and sustainable introduction of the technologies and additionally provide support and education to their population to help them to adapt to the world that will soon be upon them. With many new jobs being created, education systems will have to move away from pushing these advancements away from students to incorporating them in learning as an aid similarly to how calculators have changed the field of mathematics. As education catches up to these modern innovations – employment will start to increase and along with this comes less strain on a country’s financial systems.
\fussy